\documentclass[10pt,a4paper]{amsart}
\usepackage{amsmath,amssymb,amsthm,graphicx}
\usepackage{verbatim}
\newenvironment{reference}{\comment}{\endcomment}
%\newenvironment{reference}{}{}
\newenvironment{slogan}{\comment}{\endcomment}
\newenvironment{history}{\comment}{\endcomment}

% For commutative diagrams we use Xy-pic
\usepackage[all]{xy}

% We use 2cell for 2-commutative diagrams.
\xyoption{2cell}
% Theorem environments.
%
\theoremstyle{plain}
\newtheorem{theorem}[subsection]{Theorem}
\newtheorem{proposition}[subsection]{Proposition}
\newtheorem{lemma}[subsection]{Lemma}

\theoremstyle{definition}
\newtheorem{definition}[subsection]{Definition}
\newtheorem{example}[subsection]{Example}
\newtheorem{exercise}[subsection]{Exercise}
\newtheorem{situation}[subsection]{Situation}

\theoremstyle{remark}
\newtheorem{remark}[subsection]{Remark}
\newtheorem{remarks}[subsection]{Remarks}

\numberwithin{equation}{subsection}

% Macros
%
\def\lim{\mathop{\mathrm{lim}}\nolimits}
\def\colim{\mathop{\mathrm{colim}}\nolimits}
\def\Spec{\mathop{\mathrm{Spec}}}
\def\Hom{\mathop{\mathrm{Hom}}\nolimits}
\def\Ext{\mathop{\mathrm{Ext}}\nolimits}
\def\SheafHom{\mathop{\mathcal{H}\!\mathit{om}}\nolimits}
\def\SheafExt{\mathop{\mathcal{E}\!\mathit{xt}}\nolimits}
\def\Sch{\mathit{Sch}}
\def\Mor{\mathop{\mathrm{Mor}}\nolimits}
\def\Ob{\mathop{\mathrm{Ob}}\nolimits}
\def\Sh{\mathop{\mathit{Sh}}\nolimits}
\def\NL{\mathop{N\!L}\nolimits}
\def\CH{\mathop{\mathrm{CH}}\nolimits}
\def\proetale{{pro\text{-}\acute{e}tale}}
\def\etale{{\acute{e}tale}}
\def\QCoh{\mathit{QCoh}}
\def\Ker{\mathop{\mathrm{Ker}}}
\def\Im{\mathop{\mathrm{Im}}}
\def\Coker{\mathop{\mathrm{Coker}}}
\def\Coim{\mathop{\mathrm{Coim}}}

% Boxtimes
%
\DeclareMathSymbol{\boxtimes}{\mathbin}{AMSa}{"02}

%
% Macros for moduli stacks/spaces
%
\def\QCohstack{\mathcal{QC}\!\mathit{oh}}
\def\Cohstack{\mathcal{C}\!\mathit{oh}}
\def\Spacesstack{\mathcal{S}\!\mathit{paces}}
\def\Quotfunctor{\mathrm{Quot}}
\def\Hilbfunctor{\mathrm{Hilb}}
\def\Curvesstack{\mathcal{C}\!\mathit{urves}}
\def\Polarizedstack{\mathcal{P}\!\mathit{olarized}}
\def\Complexesstack{\mathcal{C}\!\mathit{omplexes}}
% \Pic is the operator that assigns to X its picard group, usage \Pic(X)
% \Picardstack_{X/B} denotes the Picard stack of X over B
% \Picardfunctor_{X/B} denotes the Picard functor of X over B
\def\Pic{\mathop{\mathrm{Pic}}\nolimits}
\def\Picardstack{\mathcal{P}\!\mathit{ic}}
\def\Picardfunctor{\mathrm{Pic}}
\def\Deformationcategory{\mathcal{D}\!\mathit{ef}}
\usepackage{cleveref}
\crefname{originalreference}{original reference}{original references}
\Crefname{originalreference}{Original reference}{Original references}
\makeatletter
\newlabel{stacks-sheaves-lemma-quasi-coherent-algebraic-stack}{{06WV}{source}{}{sourcepointer-1}{}}
\newlabel{stacks-sheaves-lemma-quasi-coherent-algebraic-stack@cref}{{[originalreference][0][]06WV}{[1][0][]source}{}{}{}}
\makeatother
\title{Flat pullback of quasi-coherent modules on stacks is exact}
\author{The Stacks Project Authors}
\date{stacks-project:stacks-cohomology.tex:local-snapshot-20261012; GNU FDL 1.2 or later}
\begin{document}
\maketitle
\noindent\textit{Editorial note.} All delivered mathematical text is the author's own, copied verbatim from the pinned Stacks Project chapter file stacks-cohomology.tex; no statement or proof step was rewritten, reordered or removed.\par\medskip
\noindent\textit{Editorial note.} The counted span cites the lemmas recorded in assembly.external\_references with canonical tags as navigation pointers.\par\medskip
\noindent\textit{Allowed standard prerequisites (editorial).}\begin{itemize}
\item QCoh on a presentation is abelian, pullback faithful exact. Sheaves on Stacks, Lemma stacks-sheaves-lemma-quasi-coherent-algebraic-stack (proof). Vertical functors faithful exact.
\item Flat pullback on schemes is exact. Standard scheme fact. (f')\textasciicircum{}* exact.
\end{itemize}\medskip
\ifdefined\hypertarget\hypertarget{sourcepointer-1}{}\fi
\noindent\textit{Source reference.} Original reference 06WV: sources/stacks-chow/tags.txt, source lines 8829--8829. Cross-chapter standard prerequisite cited by the author inside the counted span; it is used as a pointer and its own proof is not part of this item.\par\medskip
\begin{lemma}
\label{lemma-flat-pullback-quasi-coherent}
If $f : \mathcal{X} \to \mathcal{Y}$ is a flat morphism of algebraic stacks
then $f^* : \QCoh(\mathcal{O}_\mathcal{Y}) \to
\QCoh(\mathcal{O}_\mathcal{X})$ is an exact functor.
\end{lemma}

\begin{proof}
Choose a scheme $V$ and a surjective smooth morphism $V \to \mathcal{Y}$.
Choose a scheme $U$ and a surjective smooth morphism
$U \to V \times_\mathcal{Y} \mathcal{X}$. Then $U \to \mathcal{X}$ is
still smooth and surjective as a composition of two such morphisms.
From the commutative diagram
$$
\xymatrix{
U \ar[d] \ar[r]_{f'} & V \ar[d] \\
\mathcal{X} \ar[r]^f & \mathcal{Y}
}
$$
we obtain a commutative diagram
$$
\xymatrix{
\QCoh(\mathcal{O}_U) & \QCoh(\mathcal{O}_V) \ar[l] \\
\QCoh(\mathcal{O}_\mathcal{X}) \ar[u] &
\QCoh(\mathcal{O}_\mathcal{Y}) \ar[l] \ar[u]
}
$$
of abelian categories. Our proof that the bottom two categories in this
diagram are abelian showed that the vertical functors are faithful
exact functors (see proof of
Sheaves on Stacks, Lemma
\ref{stacks-sheaves-lemma-quasi-coherent-algebraic-stack}).
Since $f'$ is a flat morphism of schemes (by our definition of
flat morphisms of algebraic stacks) we see that $(f')^*$ is an
exact functor on quasi-coherent sheaves on $V$. Thus we win.
\end{proof}

\end{document}
